# Free-Energy-Lutar Active Inference with Hierarchical Predictive Coding and Cognitive Map Navigation

**Author:** Stephen Paul Lutar Jr.
**Affiliation:** SZL Consulting Ltd
**ORCID:** [0009-0001-0110-4173](https://orcid.org/0009-0001-0110-4173)
**Contact:** stephenlutar2@gmail.com
**Date:** 4 May 2026
**Version:** v8
**Companion papers:** v1--v7 (see References)
**Reference implementation:** [github.com/szl-holdings/szl-holdings-platform](https://github.com/szl-holdings/szl-holdings-platform) -- `packages/ouroboros-integrations/src/sovereign-engine.ts` (FELAI, PCEM, CMN classes)
**License:** CC BY 4.0

---

## Abstract

We extend Friston's Free Energy Principle (FEP) with the Lutar Omega signature (v4) to create Free-Energy-Lutar Active Inference (FELAI). FELAI replaces the standard KL divergence with an Omega-weighted divergence \( D_\Omega[q \| p] = D_{KL}[q \| p] \cdot L_\Omega \), coupling thermodynamic efficiency to the seven-version Lutar hierarchy. We integrate two neuroscience-inspired modules: (1) a Predictive Coding Error Minimizer (PCEM) implementing Rao-Ballard (1999) hierarchical prediction error minimization as a biologically plausible alternative to backpropagation, and (2) a Cognitive Map Navigator (CMN) implementing Tolman (1948) cognitive maps with O'Keefe-Moser grid cells for latent space navigation. This extends the thesis lineage from v7 (memory) to active inference and navigation.

This paper does **not** claim that FELAI outperforms standard active inference on control benchmarks. The contribution is the formal integration.

---

## 1. Motivation and Continuity from v3--v7

The v7 paper established the memory substrate (SCL, KTM, HAAM, OCM). The present paper addresses how an agent *acts* using that memory: active inference selects actions that minimize expected free energy, predictive coding generates expectations about the world, and cognitive maps provide the navigational scaffold.

The Free Energy Principle (Friston, 2010) proposes that biological systems minimize variational free energy \( F = D_{KL}[q(\theta) \| p(\theta | y)] - \ln p(y) \). Active inference extends this by treating action as a means of minimizing expected free energy. Predictive coding (Rao and Ballard, 1999) provides a neural implementation. Cognitive maps (Tolman, 1948; O'Keefe and Moser, 2014 Nobel) provide the spatial scaffold.

---

## 2. Free-Energy-Lutar Active Inference (FELAI)

### 2.1 Lutar Free Energy

\[
F_{\text{Lutar}} = L_\Omega \cdot D_{KL}[q(\theta) \| p(\theta)] + H[q(\theta)] - L_\Omega \cdot \mathbb{E}_q[\ln p(y | \theta)]
\]

where \( L_\Omega \) is the v4 Lutar Omega composite signature and \( H[\cdot] \) is entropy.

When \( L_\Omega = 1 \): FELAI reduces to standard FEP. When \( L_\Omega < 1 \): KL penalty attenuated, broader exploration. When \( L_\Omega > 1 \): penalty amplified, tighter inference.

### 2.2 Active Inference Policy Selection

Given policies \( \pi_1, \ldots, \pi_K \):

\[
\pi^* = \arg\min_{\pi_k} \sum_{t=1}^{T} \mathbb{E}_{q(\theta|\pi_k)}\left[ F_{\text{Lutar}}^{(t)} \right]
\]

Monte Carlo rollouts estimate the expected free energy over \( T \)-step horizons.

---

## 3. Predictive Coding Error Minimizer (PCEM)

### 3.1 Hierarchical Predictive Model

A predictive coding hierarchy with \( L \) layers maintains:

- Representations \( \mathbf{r}^{(l)} \) at each layer
- Top-down weight matrices \( W^{(l)} \) generating predictions
- Prediction errors \( \boldsymbol{\varepsilon}^{(l)} = \mathbf{r}^{(l)} - f(W^{(l)} \mathbf{r}^{(l+1)}) \) where \( f = \tanh \)

Total free energy across the hierarchy:

\[
F_{\text{total}} = \sum_{l=0}^{L-2} \frac{1}{2} \|\mathbf{r}^{(l)} - f(W^{(l)} \mathbf{r}^{(l+1)})\|^2
\]

### 3.2 Biologically Plausible Update Rules

Unlike backpropagation (which requires non-local weight transport), predictive coding updates are local:

\[
\Delta \mathbf{r}^{(l)} = -\eta_r \left( \boldsymbol{\varepsilon}^{(l)} - (W^{(l-1)})^T \boldsymbol{\varepsilon}^{(l-1)} \right)
\]

\[
\Delta W^{(l)} = -\eta_w \, \boldsymbol{\varepsilon}^{(l)} \, (\mathbf{r}^{(l+1)})^T
\]

The first equation adjusts representations using only local prediction errors from adjacent layers. The second is a Hebbian-like rule involving only pre-synaptic activity and post-synaptic prediction error.

**Convergence:** Under Lipschitz continuity of \( f \) and bounded weights, the predictive coding update rules converge to a local minimum of \( F_{\text{total}} \) (Millidge et al., 2021).

### 3.3 Omega-Weighted Prediction Error

\[
\bar{\varepsilon}_\Omega = \frac{1}{L-1} \sum_{l=0}^{L-2} \|\boldsymbol{\varepsilon}^{(l)}\| \cdot w_l
\]

where \( w_l \) is the v4 Omega weight for layer \( l \), assigning higher importance to prediction errors at more abstract levels.

---

## 4. Cognitive Map Navigator (CMN)

### 4.1 Cognitive Map Graph

A cognitive map \( G_{\text{cog}} = (V, E, \tau, \mathbf{p}) \) consists of:

- Nodes \( V \) representing spatial or conceptual locations
- Edges \( E \) representing navigable transitions
- Cell type function \( \tau: V \to \{\text{place, grid, head\_direction, boundary}\} \)
- Position function \( \mathbf{p}: V \to \mathbb{R}^2 \)

### 4.2 Grid Cell Firing Model

\[
g(x, y) = \frac{1}{3} \sum_{k=0}^{2} \cos\left(\frac{2\pi f}{\lambda} (x \cos\theta_k + y \sin\theta_k)\right)
\]

where \( f \) is spatial frequency, \( \lambda \) is grid spacing, and \( \theta_k = k \cdot \pi/3 \) are the three hexagonal orientations.

### 4.3 Path Integration and Spatial Coherence

Head direction signal:

\[
\text{HD}(\mathbf{p}_1, \mathbf{p}_2) = \arctan\left(\frac{p_{2,y} - p_{1,y}}{p_{2,x} - p_{1,x}}\right)
\]

Spatial coherence (ratio of direct distance to path distance):

\[
C_{\text{spatial}} = \frac{d_{\text{direct}}(\text{start}, \text{goal})}{d_{\text{path}}(\text{start}, \text{goal})}
\]

A coherent navigation achieves \( C_{\text{spatial}} = 1 \) (straight-line path).

---

## 5. Integration: FELAI + PCEM + CMN

The three modules integrate in a biologically inspired architecture:

1. **CMN** provides the latent space topology --- the map of possible states
2. **PCEM** generates predictions about the next state and computes prediction errors
3. **FELAI** uses Omega-weighted free energy to select actions minimizing expected surprise

This mirrors hippocampal cognitive maps informing cortical predictive coding, which drives active inference through the basal ganglia.

### 5.1 Connection to v7 Memory

The v7 HAAM (Hopfield associative memory) provides rapid pattern completion for CMN node lookup. KTM stores the cognitive map structure in its Core tier. OCM maintains session-local navigation context. SCL prevents catastrophic forgetting of learned maps during continual learning.

---

## 6. What this paper does and does not claim

**Claims:**
1. FELAI is a well-defined extension of the FEP with Omega weighting.
2. PCEM convergence follows from standard predictive coding theory (Millidge et al., 2021).
3. The CMN grid cell model follows the Hafting et al. (2005) parameterization.
4. All components are implemented in the public reference implementation.

**Non-claims:**
- No claim of empirical superiority over standard active inference on control tasks.
- No claim that the biological analogies have literal neuroscientific validity.
- No claim of deployed production use.

---

## 7. Conclusion

FELAI + PCEM + CMN provide a unified framework bridging active inference, predictive processing, and spatial cognition for AI systems. The integration extends the v3--v7 thesis lineage to the action-selection and navigation layers.

---

## References

1. Friston, K. (2010). The free-energy principle: a unified brain theory? *Nature Reviews Neuroscience*, 11:127--138.
2. Rao, R.P.N. and Ballard, D.H. (1999). Predictive coding in the visual cortex. *Nature Neuroscience*, 2:79--87.
3. Tolman, E.C. (1948). Cognitive maps in rats and men. *Psychological Review*, 55(4):189--208.
4. O'Keefe, J. and Dostrovsky, J. (1971). The hippocampus as a spatial map. *Brain Research*, 34:171--175.
5. Hafting, T. et al. (2005). Microstructure of a spatial map in the entorhinal cortex. *Nature*, 436:801--806.
6. Millidge, B. et al. (2021). Predictive coding approximates backprop along arbitrary computation graphs. *NeurIPS*.
7. Kirsanov, A. (2024). How Your Brain Organizes Information. YouTube lecture series.
8. Lutar, S.P. (2026a--g). Ouroboros Thesis series v1--v7. See companion papers.
